% Pessimism principle -> CQL -> BCQ/IQL. The optimism/pessimism mirror reuses
% the UCB optimism idea from exploration, flipped in sign.

\begin{frame}{The Fix: Optimism's Mirror Twin}
\definecolor{ogreen}{RGB}{106,168,79}
\definecolor{ored}{RGB}{158,28,38}
\definecolor{oblue}{RGB}{16,53,103}
\centering
\begin{tikzpicture}[font=\scriptsize, >={Stealth[length=2.4mm]}]
    \draw[very thick] (4,0.4) -- (4,2.6);
    \draw[very thick] (3.85,0.4) -- (4.15,0.4);
    \draw[very thick] (3.85,2.6) -- (4.15,2.6);
    \fill[oblue] (4,1.5) circle (2.2pt);
    \node[right=3pt, font=\scriptsize] at (4,1.5) {$\hat Q(s,a)$};
    \node[right=3pt, font=\scriptsize] at (4,2.6) {$\hat Q + u$};
    \node[left=3pt, font=\scriptsize] at (4,0.4) {$\hat Q - u$};
    \node[align=center, ogreen!60!black, font=\scriptsize\bfseries] (on) at (1.1,2.6)
      {Online (Day 7)\\optimism $\to$ explore};
    \draw[->, ogreen!70!black, thick] (on.east) -- (3.8,2.6);
    \node[align=center, ored, font=\scriptsize\bfseries] (off) at (7.1,0.4)
      {Offline (today)\\pessimism $\to$ stay supported};
    \draw[->, ored, thick] (off.west) -- (4.2,0.4);
\end{tikzpicture}
\begin{itemize}
    \item Day 7's rule: \emph{in the face of uncertainty, be optimistic}: UCB
          acts on $\hat Q + u$, because uncertainty was an \textbf{opportunity}
    \item Offline, uncertainty can never be resolved: it is a \textbf{hazard}
    \item Same machinery, opposite sign: act on $\hat Q - u$, \emph{pessimism}
          in the face of uncertainty
    \item Two places to install it: in the \textbf{values} (CQL) or in the
          \textbf{policy} (BCQ, IQL)
\end{itemize}
\end{frame}

% === CQL objective ======================================================
% The annotated objective is repeated verbatim across this build so the
% formula never jumps position between frames.
\newcommand{\cqlobjective}{%
  {\small
  \[ \min_Q\;\; \alpha\,\mathbb{E}_{s\sim\mathcal{D}}\Big[
      \underbrace{\mathbb{E}_{a\sim\pi(\cdot\mid s)}\!\left[Q(s,a)\right]}
        _{\substack{\text{what $Q$ \emph{claims} $\pi$}\\[-1pt]\text{can get (often OOD)}}}
      \;-\;
      \underbrace{\mathbb{E}_{a\sim\pi_\beta(\cdot\mid s)}\!\left[Q(s,a)\right]}
        _{\substack{\text{what $Q$ says about}\\[-1pt]\text{actions we really saw}}}
    \Big] + \mathcal{L}_{\mathrm{Bellman}} \]}}

\begin{frame}{Conservative Q-Learning (CQL)}
\begin{itemize}
    \item Train $Q$ with an extra penalty that makes unfamiliar (OOD) actions
          \emph{look bad on purpose}
    \item[] \cqlobjective
    \item $\min_Q$ = train $Q$ (tune its weights) to make this whole expression
          small. Both expectations are at the \emph{same} $s$; $\pi_\beta$ is the
          behavior policy from before --- the actions the dataset actually
          logged at $s$
\end{itemize}
\end{frame}

\begin{frame}{Conservative Q-Learning (CQL)}
\begin{itemize}
    \item Train $Q$ with an extra penalty that makes unfamiliar (OOD) actions
          \emph{look bad on purpose}
    \item[] \cqlobjective
    \item Reads as: \emph{``how much better does $Q$ claim I can do by leaving
          the data?''} --- a per-state measure of the phantom
    \item $\alpha$ = how hard to push ($\alpha{=}0$ is plain Q-learning);
          $\mathcal{L}_{\mathrm{Bellman}}$ = the normal TD loss, fitting $Q$ to
          $r + \gamma \max_{a'} Q(s',a')$
\end{itemize}
\end{frame}

\begin{frame}{Conservative Q-Learning (CQL)}
\begin{itemize}
    \item Train $Q$ with an extra penalty that makes unfamiliar (OOD) actions
          \emph{look bad on purpose}
    \item[] \cqlobjective
    \item Reads as: \emph{``how much better does $Q$ claim I can do by leaving
          the data?''} --- a per-state measure of the phantom
    \item \textbf{Why the subtraction?} Lowering $Q$ \emph{everywhere} changes no
          $\arg\max$, so it changes no policy. The gap is shift-invariant: it can
          only shrink by \emph{tilting} $Q$ down on OOD actions \emph{relative to}
          in-data ones
    \item If $\pi = \pi_\beta$ the penalty is $0$: no punishment for staying home
\end{itemize}
\end{frame}

% === What CQL does (figure) =============================================
\begin{frame}{What CQL Does to the Phantom}
\definecolor{oteal}{RGB}{0,141,160}
\definecolor{ored}{RGB}{158,28,38}
\definecolor{oblue}{RGB}{16,53,103}
\centering
\begin{tikzpicture}[font=\scriptsize, >={Stealth[length=2.4mm]}]
    \draw[->] (0,0) -- (8.0,0) node[below left]{$a$};
    \draw[->] (0,0) -- (0,3.3);
    \node[rotate=90] at (-0.35,1.6) {$Q(s,\cdot)$};
    \fill[oteal!12] (0.8,0.02) rectangle (3.6,3.0);
    \node[oteal!60!black] at (2.2,2.8) {data support};
    \draw[ored, thick, dashed] plot[smooth] coordinates
      {(0.3,1.0) (1.2,1.45) (2.2,1.95) (3.2,1.75) (4.4,1.9) (5.6,2.5) (7.0,3.05)};
    \node[ored, font=\scriptsize] at (7.35,2.85) {naive $\hat Q$};
    \draw[oblue, very thick] plot[smooth] coordinates
      {(0.3,0.95) (1.2,1.4) (2.2,1.9) (3.2,1.65) (4.4,1.05) (5.6,0.6) (7.0,0.35)};
    \node[oblue, font=\scriptsize] at (7.4,0.55) {CQL $\hat Q$};
    \foreach \x/\y in {1.0/1.42, 1.5/1.66, 1.9/1.84, 2.4/1.92, 2.8/1.85, 3.3/1.66}
      \fill[oblue] (\x,\y) circle (1.6pt);
    \draw[->, oblue, thick] (5.0,2.05) -- (5.0,0.95);
    \draw[->, oblue, thick] (6.2,2.65) -- (6.2,0.60);
    \node[oblue, font=\scriptsize, align=center] at (5.6,-0.45)
      {pushed down: $a\sim\pi$, $a\notin\mathcal{D}$};
\end{tikzpicture}
\begin{itemize}
    \item The phantom peak from before, after CQL training
    \item OOD actions now look \emph{worse} than in-data actions: the greedy
          policy stays home
\end{itemize}
\end{frame}

% === BCQ & IQL ==========================================================
\begin{frame}{A Second Family: Constrain the Policy}
\begin{itemize}
    \item CQL enforced pessimism in the \textbf{values}. A second family enforces
          it in the \textbf{policy} instead:
    \item \textbf{Keep the policy close to the data}: let it pick only actions
          the behavior policy plausibly could have taken, so it never reaches an
          OOD action in the first place
    \item \textbf{BCQ} and \textbf{IQL} are the well-known instances; they
          differ in \emph{how} they enforce this, and those details are genuinely
          involved (beyond our scope)
    \item Same principle as CQL, applied in a different place: \emph{stay where the
          data is}
\end{itemize}
\end{frame}
